\clearpage
\section{Problem Analysis}

The problem can be decomposed into four main requirements.
\begin{itemize}
    \item At most one process may enter the critical section at a time.
    \item The middleware must hide coordination details from the client code.
    \item A crash of the current owner must not leave the section permanently locked.
    \item Initialization must be safe even when multiple processes start concurrently.
\end{itemize}

The first requirement is the classical mutual-exclusion constraint, but the
other three are what make the exercise genuinely distributed. If the current
owner crashes, the system must not depend on that process to release the
critical section. Likewise, if two processes start at nearly the same time,
they must not both conclude that the first token still needs to be created.

This immediately suggests that the broker, not the client process, should be
the \textbf{single source of truth}. The middleware therefore relies on broker
queues, consumer ownership, and requeue semantics instead of maintaining a
second local copy of the distributed state.
